% Generated draft: requires mathematical and source review.
\documentclass[11pt]{article}
\usepackage[margin=25mm]{geometry}
\usepackage{fontspec}
\setmainfont{Latin Modern Roman}
\usepackage{amsmath,amssymb,mathtools,unicode-math}
\setmathfont{Latin Modern Math}
\usepackage{enumitem,needspace,xurl,hyperref}
\setlength{\parindent}{0pt}
\setlength{\parskip}{5pt}
\setlength{\emergencystretch}{3em}
\newcommand{\N}{\mathbb{N}}
\newcommand{\Z}{\mathbb{Z}}
\newcommand{\Q}{\mathbb{Q}}
\newcommand{\R}{\mathbb{R}}
\newcommand{\C}{\mathbb{C}}
\newcommand{\F}{\mathbb{F}}
\newcommand{\A}{\mathbb{A}}
\newcommand{\PP}{\mathbb{P}}
\newcommand{\E}{\mathbb{E}}
\newcommand{\eps}{\varepsilon}
\newcommand{\dd}{\,\mathrm{d}}
\begin{document}
\section*{UnsolvedMath ID: 20002714}
\textbf{Source problem code:} catalogue entry
\subsection{Definitions and mathematical statement}
\textbf{Problem 3.4 (D. Bessis).} Is there a structure theory of Lie groups or algebraic groups analogous to the dual braid monoid, and is there a dual notion of \(BN\)-pairs? The source does not specify axioms for the desired dual structure or a class of groups. The extraction's following section ``4. Garside Structures'' is neighboring exposition, not part of the numbered question; it is retained below as context.

\textbf{Reported obstruction to three naive templates.} Let \((W,S)\) be a finite Coxeter system, let \(T\) be its reflections, and let \(\ell_T\) denote reflection length. For a Coxeter element \(c\), define absolute order and the noncrossing interval by
\[
 u\leq_Tv\quad\Longleftrightarrow\quad
 \ell_T(u)+\ell_T(u^{-1}v)=\ell_T(v),\qquad
 \operatorname{NC}(W,c)=[1,c]_T.
\]
Assume some pair \(s,t\in S\) has \(m_{st}\geq3\); irreducibility is not required. The aid reports the following three obstructions.
\begin{enumerate}
\item Although \(c\in\operatorname{NC}(W,c)\), one has \(c^2\notin\operatorname{NC}(W,c)\). Hence this set is not a group under the inherited multiplication in \(W\), and cannot serve unchanged as an ordinary Weyl quotient with that multiplication. This does not forbid an unrelated group structure on a set of the same cardinality.
\item Whenever a group \(G\) has an ordinary \(BN\)-pair with Weyl group \(W=N/(B\cap N)\), and Bruhat decomposition \(G=\bigsqcup_{w\in W}B\dot wB\), the union
\[
 \bigcup_{w\in\operatorname{NC}(W,c)}B\dot wB
\]
is not a subgroup. It contains a representative of \(c\), while its square has Weyl label \(c^2\), which is excluded. This is conditional on such a \(BN\)-pair existing and does not assert that every noncrystallographic \(W\) occurs for a Lie group.
\item The poset \(\operatorname{NC}(W,c)\) contains a rank-two dihedral interval with \(m_{st}\) atoms. Every rank-two interval in Coxeter Bruhat order has exactly two atoms. Therefore \(\operatorname{NC}(W,c)\) is not isomorphic to a Coxeter Bruhat interval.
\end{enumerate}
The noncommuting-edge assumption matters: products of commuting \(A_1\) factors are outside this obstruction. The conclusions exclude retaining ordinary group multiplication, ordinary Bruhat-cell subgroup unions or ordinary Bruhat intervals as the proposed dual structure. They do not rule out Garside categories, different cells/orders or genuinely new axioms for a dual \(BN\)-pair.

\textbf{Spilled-over Garside context.} The next section credits David Bessis [6] and Tom Brady [12] with the general noncrossing lattice in their study of Artin groups \(A(W)\). It describes labeled Hasse diagrams yielding monoids \(M(P)\) and groups \(G(P)\), with appropriate interval-path relations. Additional labeling/Garside axioms give useful normal forms and word/conjugacy algorithms; those axioms are not supplied, so such algorithmic conclusions are not assigned to every arbitrary labeled poset. Jon McCammond's survey [35] is recommended.

The aid reports later analogues involving Garside categories, Deligne--Lusztig theory, complex reflection groups and categorical ADE constructions, but says it found no accepted \(\operatorname{NC}(W,c)\)-indexed dual-\(BN\) structure theorem in the literature checked. No full references, date or precise intended axioms accompany that assessment. It is not a nonexistence result, and the three elementary template obstructions do not settle the open-ended source question.
\subsection{Short English statement}
Bessis asks for a dual-braid analogue of group structure theory and \(BN\)-pairs. The aid excludes three direct templates when the Coxeter system has a noncommuting edge: inherited noncrossing multiplication, a subgroup union of ordinary Bruhat cells, and an ordinary Bruhat interval. The adjacent Garside exposition and broader possible analogues are retained separately.
\subsection{Sources}
\begin{itemize}
\item Local UnsolvedMath registry, permanent ID \(20002714\), catalogue entry, \texttt{aim-probability-notes.json}, \texttt{notes[155]}, workshop \emph{Braid groups, clusters and free probability}, Problem 3.4, D. Bessis, and appended Section 4: \url{https://aimath.org/WWN/braidgroups/braidgroups.pdf}.
\item Original references by label: David Bessis [6], dual braid monoids; Tom Brady [12], noncrossing partitions and Artin groups; Jon McCammond [35], Garside survey. Full bibliography entries are not supplied.
\item Supplied research summary on the three template obstructions and later Garside-category/Deligne--Lusztig/complex-reflection/categorical ADE analogues. Attempt 1: partial\_result, status at run ``open'', rigor ``proved'', low novelty confidence; registry status partially\_solved. The machine-generated aid is not peer reviewed and requires independent verification before citation. The supplied verification reports page 8 inspection; this draft has not visited the PDF.
\end{itemize}
\end{document}
